\section{Discussion}\label{sec:discussion}

This paper develops a common conditioning construction for positivity,
monotonicity, and convexity in the Laplacian eigenfunction expansion.
Its contribution is to express these whole-interval events within
the HSGP model using its existing basis and kernel-derived Gaussian
prior, including joint restrictions when a suitable common direction
is available. The resulting representation supports joint inference
for the curve and its derivatives and separates the definition of
the conditioned posterior from the numerical evaluation of its
threshold.

The approximation theory connects this finite construction to
conditioning the underlying GP. For the individual constraints,
the connection holds under the stated spectral assumptions as both
padding and frequency cutoff increase, giving a basis for using
finite expansions to approximate the corresponding constrained
process. The joint construction is exact within the finite expansion,
while the convergence theorem covers the individual restrictions.

The practical cost of the representation comes from evaluating a
global threshold at each posterior-density calculation. Its
maximizing location can change with the coefficients, creating kinks
that matter for gradient-based sampling even when the marginal law
is strongly log-concave. The threshold-error analysis controls the
posterior change caused by conservative evaluation, while sampler
accuracy and computational cost remain separate considerations.
Developing samplers that exploit this geometry while accommodating
the nonsmooth threshold is a useful direction for further work.

The mercury and leaf applications show that shape restrictions can
substantially change derivative inference even when the fitted response
curves remain similar. These changes matter because the derivatives
determine the inferred mercury concentration and marginal carbon
assimilation, respectively. Comparing the constrained and ordinary
Gaussian posteriors makes the contribution of the shape assumption
explicit in both the estimates and their uncertainty. The construction
thus connects scientifically motivated restrictions to joint inference
for the response and its derivatives, with the chosen constraint
enforced throughout the interval.
